\section{Method}
\label{sec:method}


In this section,
we introduce our Vlogger framework in detail.
First,
we describe how to make a vlog by planning and shooting with our Vlogger.
Then,
we further explain the novel design of ShowMaker.
As the videographer in our Vlogger,
it is critical for video generation in the vlog.


\subsection{Overall Framework of Vlogger}

To generate a minute-level vlog,
our Vlogger leverages LLM as the director,
which can effectively decompose this generation task by four key roles within the planning and shooting stages. 
As shown in Fig. \ref{fig:teaser},
the LLM Director first creates Script and designs Actor in the planning stage.
Based on Script and Actor,
ShowMaker generates a video snippet for each scene in the shooting stage,
and Voicer dubs subtitles of this snippet.
Finally,
one can combine the snippets of all the scenes as a vlog.


\begin{figure*}[t]
    \centering
    % \vspace{-0.3cm}
    \includegraphics[width=0.95\linewidth]{figs/generating.pdf}
    \vspace{-0.2cm}
    \caption{
    \textbf{Bottom-Up Shooting.} 
    For each scene,
    ShowMaker can generate the video snippet with script and actor coherence,
    by using script description and actor image as textual and visual prompts.
    Moreover,
    ShowMaker can effectively control the snippet duration,
    by performing generation and prediction sequentially in the inference.  
    Finally,
    we apply a Text-To-Speech (TTS) model as Voicer for dubbing.
    }
    \label{fig:generating}
    \vspace{-0.3cm}
\end{figure*}


\subsubsection{Top-Down Planning}

A vlog often comes from a user story that contains diversified content within many shot changes.
Clearly,
it is challenging to produce a minute-long vlog,
by directly feeding such a complex story into video generation models.
To this end,
we propose to leverage LLM as director and decompose the user story by top-down planning in Fig. \ref{fig:director}.

\textbf{Script Creating}.
First,
we parse the use story into a script,
which describes this story explicitly by a number of shooting scenes.
In this case,
we can generate a video snippet for each shooting scene,
instead of learning a long video tediously from the entire story.
Since LLM has shown an impressive capacity in language understanding \cite{brown2020language, Zeng2022GLM130BAO, Touvron2023LLaMAOA, OpenAI2023GPT4TR, Sun2021ERNIE3L, Bai2023QwenTR, 2023internlm, Chowdhery2022PaLMSL},
we feed the user story into such a director for script generation.
As shown in Fig. \ref{fig:director},
we introduce a progressive creation paradigm,
which can effectively parse the story by coarse-to-fine steps,
\begin{equation}
\mathcal{S}^{(i)}=\text{LLM}(\mathcal{S}^{(i-1)},\mathcal{I}^{(i)}_{S},\mathcal{U}).
\label{eq:script}    
\end{equation}
Given the user story $\mathcal{U}$ and creation instruction $\mathcal{I}^{(i)}_{S}$,
the LLM Director generates the current script $\mathcal{S}^{(i)}$ from the previous one $\mathcal{S}^{(i-1)}$.
More specifically,
there are four steps including
(1) \textit{Rough}.
First,
LLM generates a basic draft of the script from the story.
(2) \textit{Detailed}.
Then,
LLM refines the rough script with story details.
(3) \textit{Completed}.
Next,
LLM checks if the detailed script misses the important parts of the story. 
(4) \textit{Scheduled}.
Finally,
LLM allocates a shooting duration for each scene in the completed script,
according to the scene content.
For convenience,
we denote the final script as $\mathcal{S}$ in the following.
It contains the descriptions of $N$ shooting scenes $\{\mathcal{S}_{1},...,\mathcal{S}_{N}\}$ and their allocated duration $\{\mathcal{T}_{1},...,\mathcal{T}_{N}\}$.
Additionally,
due to the limited pages,
please read the full descriptions of instructions and script in the supplementary doc.


\textbf{Actor Designing}.
After generating the script,
it is time to design actors in the vlog.
As shown in Fig. \ref{fig:director},
we ask the LLM Director to summarize the actor list $\mathcal{A}$ from script $\mathcal{S}$, %and describe appearance of these actors $\mathcal{C}$
\begin{equation}
\mathcal{A}=\text{LLM}(\mathcal{S},\mathcal{I}_{A}),
\label{eq:actor_des}    
\end{equation}
where 
$\mathcal{I}_{A}$ is the instruction for actor summarization.
Then,
according to actor descriptions,
the LLM Director invokes a designer to generate reference images of these actors $\mathcal{R}$,
\begin{equation}
\mathcal{R}=\text{Stable-Diffusion}(\mathcal{A}).
\label{eq:actor}    
\end{equation}
We choose Stable Diffusion XL \cite{Podell2023SDXLIL} as the character designer,
due to its high-quality generation.
Finally,
based on script $\mathcal{S}$ and actor $\mathcal{A}$,
the LLM Director decides the leading actor (\textit{i.e.}, protagonist) in each shooting scene of the script, 
\begin{equation}
\mathcal{P}=\text{LLM}(\mathcal{S},\mathcal{A},\mathcal{I}_{P}).
\label{eq:actorp}    
\end{equation}
where 
$\mathcal{I}_{P}$ is the instruction for protagonist selection.
The resulting doc $\mathcal{P}$ is aligned with the script $\mathcal{S}$,
where 
$\mathcal{P}_{n}$ describes which actor appears in the scene $\mathcal{S}_{n}$.
For example,
\{$\mathcal{S}_{6}$ \textit{is starred by} $\mathcal{A}_{1}$\} in Fig. \ref{fig:director}.



\begin{figure*}[t]
    \centering
    % \vspace{-0.3cm}
    \includegraphics[width=\textwidth]{figs/arch.pdf}
    %\vspace{-0.7cm}
    \caption{
    \textbf{ShowMaker}. 
    (a) The Overall Architecture.
    (b) Spatial-Temporal Enhanced Block (STEB).
    Via spatial-actor and temporal-text cross attention, 
    our STEB can further enhance actor and script coherence in the snippet.
    (c) Mode Selection. 
    We introduce a mixed training paradigm of T2V generation and prediction,
    via probabilistic selection of masked frames.
    %The {\color{gray}{gray}} modules in STEB represent the original components that are inherent to the base model.
    }
    \label{fig:architecture}
    \vspace{-0.3cm}
\end{figure*}

\subsubsection{Bottom-Up Shooting}
\label{sec:Bottom-Up}

Via the top-down planning above,
the LLM Director flexibly decomposes a complex user story into several script scenes 
and designs actor reference image for each scene.
Such a manner largely reduces the difficulty of vlog generation,
since we can generate vlogs by bottom-up shooting,
\textit{i.e.},
we just need to generate the video snippet for each shooting scene and combine all of them as a vlog.

\textbf{ShowMaker Shooting}.
To generate the video snippet of a shooting scene,
we introduce a novel ShowMaker as the videographer,
which is a video diffusion model with two distinct designs.
First,
it is important to maintain spatial-temporal coherence of both script and actor in the generated snippet.
Hence,
our ShowMaker not only takes the scene description $\mathcal{S}_{n}$ as a textual prompt 
but also takes actor image of this scene $\mathcal{R}_{n}$ as a visual prompt.
Second,
each scene is allocated with a shooting duration in the script.
To control the duration of each scene,
our ShowMaker contains two learning modes including generation and prediction in Fig.\ref{fig:generating}.
Specifically,
it starts with the generation mode.
For the shooting scene $n$,
we feed its script description $\mathcal{S}_{n}$ and actor reference image $\mathcal{R}_{n}$ into ShowMaker,
\begin{equation}
\mathcal{C}_{n}^{(1)}=\text{ShowMaker}(\mathcal{N}_{n}^{(1)}~|~\mathcal{S}_{n},\mathcal{R}_{n},\text{Generate}),
\label{eq:showmakergeneration}    
\end{equation}
which generates the first video clip of this scene $\mathcal{C}_{n}^{(1)}$ from the noisy clip $\mathcal{N}_{n}^{(1)}$.
If the duration of this clip is smaller than the allocated duration $\mathcal{T}_{n}$ in the script,
we continue to perform the prediction mode, %in an auto-aggressive manner,
\textit{i.e.},
the last $k$ frames of the current clip $\mathcal{C}_{n}^{(j)}(k)$ are used as context,
when generating the next clip from the noisy input $\mathcal{N}_{n}^{(j+1)}$,
\begin{equation}
\mathcal{C}_{n}^{(j+1)}=\text{ShowMaker}(\mathcal{N}_{n}^{(j+1)}~|~\mathcal{S}_{n},\mathcal{C}_{n}^{(j)}(k),\text{Predict}).
\label{eq:showmakerprediction}    
\end{equation}
Note that,
actor reference images are not necessary in the prediction mode,
since such an actor's appearance has been shown in the current clip $\mathcal{C}_{n}^{(j)}(k)$ that is used as input for prediction.
This prediction procedure stops until the total duration achieves the allocated $\mathcal{T}_{n}$ of this scene.
Subsequently,
we combine all the clips as the video snippet of this scene, \textit{i.e.}, $\mathcal{C}_{n}=\{\mathcal{C}_{n}^{(1)},...,\mathcal{C}_{n}^{(J)}\}$.



\textbf{Voicer Speaking}.
To enhance the completeness of the vlog,
we apply a Text-To-Speech model (e.g., Bark \cite{sunoai2023bark}) as a Voicer,
which converts the scene description $\mathcal{S}_{n}$ into the corresponding audio $\mathcal{O}_{n}=\text{Bark}(\mathcal{S}_{n})$.
Finally,
we add this audio to the corresponding video snippet
$\mathcal{V}_{n}=\mathcal{O}_{n}\oplus\mathcal{C}_{n}$,
and
combine all the sounded video snippets as a complete vlog,
\textit{i.e.},
$\mathcal{V}=\{\mathcal{V}_{1},...,\mathcal{V}_{N}\}$.



\subsection{ShowMaker}

As discussed in Section \ref{sec:Bottom-Up},
ShowMaker plays a critical role in generating video snippet of a shooting scene.
In this work,
we introduce a text-to-video diffusion model for it.
It follows the style of latent diffusion model \cite{rombach2022highResolution}.
In the diffusing stage,
we add Gaussian noise progressively to the latent code of a training snippet.
In the denoising stage,
we reconstruct the latent code from the noisy latent code at any iteration step.
For simplicity,
we just show the denoising stage in Fig.\ref{fig:architecture} (a).
First,
we forward a noisy training snippet into the encoder of the autoencoder to extract its latent code.
Then,
we feed this into a denoising U-Net \cite{Ronneberger2015UNetCN} to learn the clean latent code.
Finally,
we leverage the decoder to reconstruct the original snippet with the clean latent code.

But, compared with the existing video diffusion models \cite{Ho2022VideoDM, blattmann2023align, Wang2023LAVIEHV, Ho2022ImagenVH, Ge2023PreserveYO},
our novel ShowMaker contains two distinct designs,
in terms of 
model structure (\textit{i.e.}, Spatial-Temporal Enhanced Block) in  Fig.\ref{fig:architecture} (b),
and
training paradigm (\textit{i.e.}, Mode Selection) in Fig.\ref{fig:architecture} (c).
 

\subsubsection{Spatial-Temporal Enhanced Block (STEB)}
\label{sec:STEB}

To reconstruct the clean latent code of a training video snippet,
each block in the denoising U-Net consists of both spatial and temporal operations in previous works \cite{Wang2023LAVIEHV}.
First,
spatial operations encode the feature of each frame separately in the snippet.
Typically,
three operations are inherited from text-to-image generation approaches \cite{saharia2022imagen, Ramesh2022HierarchicalTI, rombach2022highResolution, Balaji2022eDiffITD, Nichol2021GLIDETP},
including 
spatial ConVolution (CV), 
spatial Self Attention (SA),
and 
spatial Cross Attention (CA),
\begin{align}
\mathcal{X}_{cv}&=\text{CV-Spatial}(\mathcal{X}_{in}, \mathcal{F}_{t}),
\label{eq:steb1}\\    
\mathcal{X}_{sa}&=\text{SA-Spatial}(\mathcal{X}_{cv}),
\label{eq:steb2}\\
\mathcal{X}_{ca}&=\text{CA-Spatial-Text}(\mathcal{X}_{sa},\mathcal{S}_{n}),
\label{eq:steb3} 
\end{align}
where
$\mathcal{X}_{in}$ is the noisy feature of the training snippet,
and $\mathcal{F}_{t}$ is the positional embedding of the iteration step $t$.
To guide denoising with the given text, 
we use the scene description $\mathcal{S}_{n}$ as Key and Value of cross attention in Eq. (\ref{eq:steb3}).

However,
we notice that 
such spatial encoding does not consider actors.
Hence,
it inevitably suffers from actor incoherence when generating a snippet.
To tackle this problem,
we introduce a spatial image cross attention,
\begin{equation}
\mathcal{Y}_{ca}=\text{CA-Spatial-Actor}(\mathcal{X}_{sa},\mathcal{R}_{n}).
\label{eq:steb4} 
\end{equation}
For a shooting scene $\mathcal{S}_{n}$,
we leverage the protagonist doc (Eq. \ref{eq:actorp}) to find the corresponding actor in this scene.
Then,
we leverage actor reference image $\mathcal{R}_{n}$ as the visual context of spatial cross attention.
Subsequently,
we enhance spatial embedding with the complementary guidance of both script and actor,
\textit{i.e.},
$\mathcal{Z}_{se}=\mathcal{X}_{ca}+\beta\mathcal{Y}_{ca}$ with a weight $\beta$.

After spatial encoding,
it is time to learn the correlation across frames in the snippet.
Hence,
the typical operation is to perform self attention along the temporal dimension,
\begin{equation}
\mathcal{Z}_{sa}=\text{SA-Temporal}(\mathcal{Z}_{se}).
\label{eq:steb5} 
\end{equation}
However,
such a temporal operation does not take the constraints of scene text into account.
It often leads to text incoherence when generating video snippets.
Hence,
we introduce a temporal text cross attention,
\begin{equation}
\mathcal{Z}_{ca}=\text{CA-Temporal-Text}(\mathcal{Z}_{sa},\mathcal{S}_{n}),
\label{eq:steb6} 
\end{equation}
where
we leverage the scene description $\mathcal{S}_{n}$ as the textual context of temporal encoding.
Via spatial-actor (Eq. \ref{eq:steb4}) and temporal-text (Eq. \ref{eq:steb6}) cross attention,
our STEB can further enhance actor and script coherence in the snippet.



\subsubsection{Mixed Training Paradigm with Mode Selection}

As discussed in Section \ref{sec:Bottom-Up},
ShowMaker aims to generate a video snippet with the allocated duration in the script.
To this goal,
it leverages the combination of both generation and prediction modes in the inference stage.
Next,
we explain how to train our ShowMaker to learn these modes.

\begin{table}[tp]
\centering
\begin{minipage}[t]{0.49\linewidth}
    \vspace{0pt}
    \centering
    \setlength\tabcolsep{7.0pt}
    \resizebox{0.986\textwidth}{!}{
        \begin{tabular}{l c}
        \toprule
        \textbf{Method} & \textbf{FVD ()}  \\
        \midrule        
        VideoFactory~\cite{Wang2023VideoFactorySA} & 410.00 \\
        % Show-1~\cite{Zhang2023Show1MP} & \cmark & 394.46 \\
        Make-A-Video~\cite{singer2023makeavideo} & 367.23 \\
        PYoCo~\cite{Ge2023PreserveYO} & 355.19 \\
        %% base model~\cite{base model} & \cmark & 526.30 \\
        \midrule
        \textbf{Ours} & \textbf{292.43} \\
        \bottomrule
        % \vspace{-1.8em}
    \end{tabular}
    }
    \vspace{0.2cm}
    \subcaption{Hand-crafted prompt.}
    \label{tab:ucf_manual}
\end{minipage}
%
\begin{minipage}[t]{0.49\linewidth}
    \vspace{0pt}
    \centering
    \centering
    \setlength\tabcolsep{3.8pt}
    \resizebox{1.04\textwidth}{!}{
        \begin{tabular}{l c}
        \toprule
        \textbf{Method} & \textbf{FVD ()}  \\
        \midrule
        CogVideo (Chinese)~\cite{hong2023cogvideo} & 751.34 \\
        CogVideo (English)~\cite{hong2023cogvideo} & 701.59 \\
        MagicVideo~\cite{Zhou2022MagicVideoEV} & 699.00 \\
        Video LDM~\cite{blattmann2023align} & 550.61 \\
        % LaVie~\cite{wang2023lavie} & \cmark & 526.30 \\
        %% base model~\cite{base model} & \cmark & 526.30 \\
        \midrule
        \textbf{Ours} & \textbf{525.01} \\
        \bottomrule
        % \vspace{-2em}
    \end{tabular}
    }
    \vspace{-0.1cm}
    \subcaption{Class label.}
    \label{tab: ucf_class}
\end{minipage}
\vspace{-0.3cm}
\caption{\textbf{Zero-shot comperison with the state-of-the-art methods on UCF-101}.
Hand-crafted prompt from \cite{Ge2023PreserveYO}.
}
\vspace{-0.3cm}
\label{tab:ucf}
\end{table}

As shown in Fig.\ref{fig:architecture} (c),
we design a mode selection mechanism,
which selects $k$ frames of the clean snippet as the context of the noisy snippet.
The $k$$=$$0$ setting refers to the generation mode 
since there is no context from the clean snippet.
Alternatively,
the $k$$>$$0$ setting refers to the prediction mode,
since $k$ frames of the clean snippet are already available.
The goal is to generate the rest frames of the clean snippet from the noisy one.
To integrate both modes into training,
we design a probabilistic manner to select $k$,
\begin{align}
\begin{split}
\text{P}(k)=\left\{
\begin{aligned}
\alpha^{k}-\alpha^{k+1},~~~~~~& k\in\lbrack0,m) \\
\alpha^{k},~~~~~~& k=m
\end{aligned}
\right.
\end{split}
\label{eq:alpham} 
\end{align}
where 
$\text{P}(k)$ is the selection probability distribution of $k$.
Moreover,
$0$$<$$\alpha$$<$$1$ and $0$$\leq$$m$ are the manual parameters,
which respectively control the mode selection tendencies of $\text{P}(k)$ and the maximum number of preserved frames. 

After determining $k$,
we introduce a frame mask $\mathcal{M}_{k}$ on the latent code of the clean snippet $\mathcal{X}_{clean}$,
\begin{equation}
\mathcal{X}_{k}=\mathcal{X}_{clean}\odot\mathcal{M}_{k}, \label{eq:ms1}
\end{equation}
where
we only preserve $k$ frames and mask the rest of the frames.
Then,
we concatenate the mask $\mathcal{M}_{k}$ with the latent code of snippet $\mathcal{X}_{noise}$ and $\mathcal{X}_{k}$,
\begin{equation}
\mathcal{X}_{in}=\text{Concat}(\mathcal{X}_{noise},~\mathcal{X}_{k},~\mathcal{M}_{k}).\label{eq:ms3}  
\end{equation}
This produces the input feature $\mathcal{X}_{in}$ for training the denoising U-Net in Section \ref{sec:STEB}.
Via such a concise probabilistic manner,
we can effectively integrate both generation ($k$$=$$0$) and prediction ($k$$>$$0$) modes in the training procedure.





\begin{table}[t]
    \centering
    \setlength\tabcolsep{9pt}
    \resizebox{1\linewidth}{!}{%
    \begin{tabular}{l  c c c}
        \toprule
        \textbf{Method} &  \textbf{Zero-Shot} & \textbf{Pre-training Videos} & \textbf{FID ()}  \\
        \midrule        
        T2V~\cite{li2017video} & \xmark & \xmark & 82.13 \\
        SC~\cite{balaji2019Conditional} & \xmark & \xmark & 33.51 \\
        TFGAN~\cite{balaji2019Conditional} & \xmark & \xmark & 31.76 \\
        NUWA~\cite{wu2022nwa} & \xmark & 0.97M & 28.46 \\
        \midrule
        Phenaki~\cite{villegas2023phenaki} & \cmark & 15M & 37.74 \\
        \textbf{Ours} & \cmark & 10M & \textbf{37.23} \\
        \bottomrule
        \vspace{-2em}
    \end{tabular}
    }
    \caption{\textbf{Comperison with the state-of-the-art methods on Kinetics-400}. 
Both under a zero-shot setting, our FID is better than that of Phenaki, 
with only 66.7\% pre-training videos.}
    \vspace{-1em}
    \label{tab:k400}
\end{table}
\begin{table}[t]
    \centering 
    \setlength\tabcolsep{3pt}
    \resizebox{1.0\linewidth}{!}{%
    \begin{tabular}{l  c c c}
        \toprule
        \textbf{Method} &  \textbf{Zero-Shot} & \textbf{CLIPSIM ($\uparrow$)} & \textbf{CLIP-FID ()}  \\
        \midrule
        GODIVA~\cite{Wu2021GODIVAGO} & \xmark & 0.2402 & - \\
        N\"{U}WA~\cite{wu2022nwa} & \xmark & 0.2439 & 47.68 \\
        CogVideo (Chinese)~\cite{hong2023cogvideo} & \cmark & 0.2614 & 24.78 \\
        CogVideo (English)~\cite{hong2023cogvideo} & \cmark & 0.2631 & 23.59 \\
        Video LDM~\cite{blattmann2023align} & \cmark & 0.2929 & - \\
        % ModelScopeT2V~\cite{Wang2023ModelScopeTT} & \cmark & 0.2930 & 11.09 \\
        % LaVie~\cite{wang2023lavie} & \cmark & 0.2949 & - \\
        Make-A-Video~\cite{singer2023makeavideo} & \cmark & 0.3049 & 13.17 \\
        % Show-1~\cite{Zhang2023Show1MP} & \cmark & \textbf{0.3072} & 13.08 \\
        \color{gray}{PYoCo}~\cite{Ge2023PreserveYO} & \color{gray}{\cmark} & \color{gray}{-} & \color{gray}{10.21} \\
        \midrule
        \textbf{Ours} & \cmark & 0.2908 & \textbf{12.67} \\
        \bottomrule
        \vspace{-2em}
    \end{tabular}
    }
    \caption{\textbf{Comperison with the state-of-the-art methods on MSR-VTT}. PYoCo \cite{Ge2023PreserveYO} sample 59,794 videos for CLIP-FID (noted in {\color{gray}{gray}}) while others only sample less than 3,000 videos.}
    \vspace{-1.2em}
    \label{tab:msrvtt}
\end{table}